% !TeX encoding = UTF-8

\documentclass[
  uplatex,        % エンジン指定
  dvipdfmx,       % ドライバ指定
  report,         % 章(chapter)を使えるようにする
  fontsize=11pt,  % フォントサイズ
  paper=a4,       % 用紙サイズ
  openany         % 空白ページを入れない
]{jlreq} % ドキュメントクラスはjlreq


\usepackage{graphicx}
% \usepackage{xcolor}
\usepackage{amsthm, bbm, bm, booktabs, float, listings, mathtools, setspace, siunitx, amssymb} 
\usepackage{newtxtext, newtxmath} % フォントをTimes系に
\usepackage[top=30mm,bottom=30mm,left=25mm,right=25mm]{geometry} %上下左右の余白
\usepackage[dvipdfmx,bookmarks=true,bookmarksnumbered=true,colorlinks=true,linkcolor=black,citecolor=black,urlcolor=blue]{hyperref}
\usepackage{pxjahyper} % 日本語文字化け防止
\setlength{\marginparwidth}{2cm}
\usepackage[colorinlistoftodos]{todonotes} % 付箋：最終版提出時はdisableオプションをつければ消える.使用例: \todo{ここを見直す}, \todo[inline]{本文中に大きく表示}
\usepackage[nameinlink,noabbrev]{cleveref} % スマート参照

% 定理環境の定義（番号を節ごとに振る：定理 2.1 など）
\newtheorem{theorem}{定理}[section]
\newtheorem{definition}[theorem]{定義}
\newtheorem{lemma}[theorem]{補題}
\newtheorem{proposition}[theorem]{命題}
\newtheorem{corollary}[theorem]{系}
\newtheorem{note}[theorem]{注意}
\renewcommand{\proofname}{\textbf{証明}} % 証明環境カスタマイズ（証明の終わりの□を表示）

% 日本語の接頭辞設定
\crefname{equation}{式}{式}
\crefname{figure}{図}{図}
\crefname{table}{表}{表}
\crefname{chapter}{第}{章}
\crefname{section}{}{節}
\crefname{theorem}{定理}{定理}
\crefname{definition}{定義}{定義}
\crefname{lemma}{補題}{補題}
\crefname{proposition}{命題}{命題}
\crefname{corollary}{系}{系}
\crefname{note}{注意}{注意}
\newcommand{\crefrangeconjunction}{--}

% report系の既定（chapter.section）ではなく、sectionを最大単位として採番する
\renewcommand{\thesection}{\arabic{section}}
\renewcommand{\thesubsection}{\thesection.\arabic{subsection}}
\renewcommand{\thesubsubsection}{\thesubsection.\arabic{subsubsection}}

% 数式コマンドのショートカット自作
\newcommand{\E}{\mathbb{E}} % 期待値
\newcommand{\pb}{\mathbb{P}} % 確率
\newcommand{\R}{\mathbb{R}} % 実数
\newcommand{\N}{\mathbb{N}} % 自然数
\newcommand{\var}{\operatorname{Var}} % 分散
\newcommand{\cov}{\operatorname{Cov}} % 共分散
\newcommand{\corr}{\operatorname{Corr}} % 相関係数
\newcommand{\ind}{\mathbf{1}} % 指示関数(太字).  
% \newcommand{\ind}{\mathbbm{1}} % 指示関数(白抜き). こちらはローカル環境で動く
\newcommand{\pto}{\xrightarrow{\ \ p\ \ }} % 確率収束
\newcommand{\dto}{\xrightarrow{\ \ d\ \ }} % 分布収束
\newcommand{\asto}{\xrightarrow{\ \ a.s.\ \ }} % 概収束
\newcommand{\iid}{\stackrel{\mathrm{i.i.d.}}{\sim}} % i.i.d.
\newcommand{\Nd}{\mathcal{N}} % 正規分布
\newcommand{\pdif}[2]{\frac{\partial #1}{\partial #2}} % 偏微分を簡単に書くコマンド.例: \pdif{f}{x} -> ∂f/∂x
\DeclareMathOperator*{\argmax}{arg\,max} % argmax, argmin の定義
\DeclareMathOperator*{\argmin}{arg\,min}


\numberwithin{equation}{section} % 数式番号をセクション単位にする.


% 表紙のカスタマイズ設定（ここから基本書き換えない）
\makeatletter
\newcommand{\academicYear}[1]{\def\@academicYear{#1}}
\newcommand{\advisor}[1]{\def\@advisor{#1}}
\newcommand{\studentID}[1]{\def\@studentID{#1}}
\newcommand{\affiliation}[1]{\def\@affiliation{#1}}

\renewcommand{\maketitle}{
    \begin{titlepage}
        \centering
        \vspace*{20mm}

        {\Large \@academicYear \quad 修士論文}
        \vspace*{40mm}

        \begin{spacing}{1.2}
            {\huge \bfseries \@title \par}
        \end{spacing}
        \vspace{40mm}
        
        {\Large \@date}
        \vspace{20mm}
        
        {\Large 指導教員 \quad \@advisor}
        \vspace{30mm}
        
        {\Large \@affiliation}
        \vspace{10mm}
        
        {\Large \@studentID \quad \@author}
    \end{titlepage}
}
\makeatother

% abstract環境の設定
\renewenvironment{abstract}{
    \clearpage
    \setcounter{page}{1}
    \begin{center}
        {\Large \bfseries 概要}
    \end{center}
    \vspace{10mm}
    \begin{quote}
}{
    \end{quote}
}

% 参考文献(thebibliography環境)の見出しをChapterサイズする設定
\makeatletter
\renewenvironment{thebibliography}[1]{%
    \section*{参考文献} % 強制的にSection*（番号なし節）にする
    \addcontentsline{toc}{section}{参考文献} % 目次にも「参考文献」を追加する
  \list{\@biblabel{\@arabic\c@enumiv}}%
       {\settowidth\labelwidth{\@biblabel{#1}}%
        \leftmargin\labelwidth
        \advance\leftmargin\labelsep
        \@openbib@code
        \usecounter{enumiv}%
        \let\p@enumiv\@empty
        \renewcommand\theenumiv{\@arabic\c@enumiv}}%
  \sloppy
  \clubpenalty4000
  \@clubpenalty \clubpenalty
  \widowpenalty4000%
  \sfcode`\.\@m}
 {\def\@noitemerr
   {\@latex@warning{Empty `thebibliography' environment}}%
  \endlist}
\makeatother
% --- ここまで設定なので基本書き換えない ---


% --- ここから書く ---
\begin{document}

% 表紙情報の入力（ここを書き換える）
\academicYear{令和8年度}
\title{論文テンプレート\\（ここにタイトル）}
\date{2027年2月1日}
\advisor{清水 泰隆}
\affiliation{早稲田大学 基幹理工学研究科 数学応用数理専攻}
\studentID{5125A028}
\author{中川 浩輔}

\maketitle %表紙出力

% --- 概要 --- 
\begin{abstract}
    ここに論文の概要（アブスト）を記述する．
    通常このセクションでは研究の目的，手法，主な結果，結論を簡潔にまとめる．
\end{abstract}

\pagenumbering{arabic} % ページ番号を 1, 2... にリセット
\setcounter{page}{1}   % 強制的に1ページ目にする



% --- 序論 ---
\begingroup
    \let\clearpage\relax        % 改ページ命令を無効化
    \let\cleardoublepage\relax  % 偶数奇数ページ調整も無効化
    \section{序論}
\endgroup
ここに序論を書く．この研究の背景や重要性，既存の手法について説明する．そしてこの論文で何をやったのかを説明をする．
本文で各内容を明確にするため，適当に流さず「目的（何を解決するためにこの論文を書いた？）」や「動機（なぜこの論文を書きたいと思った？）」，「意義（この論文で誰がハッピーになる？）」などを
数式をなるべく用いずに，わかりやすく簡潔に示す大事な章である．

読者（審査員）はここを読んで「面白そうか」「読む価値があるか」を判断する．
% --- 背景 ---
\section{背景}
本文を読み進めるうえで必要となる事前知識や，この研究の先行研究について説明する．
この研究に強く関わるものはより詳しく説明する．
\begin{itemize}
    \item 単なる知識の羅列（教科書の写し）にしない．提案手法で使う数式や定義を厳選して定義する．使わない定義は書かない．
    \item 確率変数$X$は大文字，実現値$x$は小文字，ベクトル$\bm{x}$は太字など，記号のルールをここで確定させ，論文を通して一貫する．
    \item 既存手法の限界を突く．「手法Aは画期的だが，〇〇の条件下では動かない」→「だから私の研究が必要」という流れを作る．
\end{itemize}

\subsection{小節}
小節2.1を書く．

\subsubsection{小々節}
小節2.1.1を書く．こまめに節や小節に分け，見やすさを意識しよう．
% --- 手法 ---
\section{手法}
論文の提案内容（新しい手法や発見など）の詳しい説明．論文の核心部分であり，あなたのオリジナリティをアピールする章である．
具体的な解決策（How）を記述し，最もページ数を割くべきパートである．

\begin{itemize}
    \item いきなり数式に入ると読者は迷子になるので，まずは直感的な理解を促すといいかも．
    \item 何を解く問題なのかをきちんと定義する．（例：「対数尤度関数 $L(\theta)$ を最大化するパラメータ $\theta$ を求めることが目的である．」）
    \item 提案手法が単に動くだけでなく，数学的に良い性質（一致性，計算オーダーなど）を持っていることも示したいところ．
    \item 計算過程を飛ばしすぎない．数式展開は丁寧に．また，図解を入れる．
\end{itemize}
% --- 評価実験 ---
\section{評価実験}
\begin{itemize}
    \item 再現性を意識する．読者がこの章を読めばプログラムをゼロから実装できるレベルで詳しく書く．（ソースコードは論文に直貼りせず，GitHub等にアップロードしてリンクを掲載するほうがいいかも．）
    \item 自分の手法だけ載せるのではなく，既存手法と比較して優れていることを示す．
\end{itemize}

\subsection{実験}

\subsubsection{実験器具}
ここに実験環境などを書く．使用したデータセット，ハイパーパラメータ，計算機環境（CPU, GPU等），評価指標などを詳細に．

\subsubsection{実験手法}
ここに実験の手順を具体的に書く．

\subsection{結果}
結果を定量的に図や表にまとめる．表やグラフを見て分かる事実を述べる（例：「手法Aより手法Bの方が精度が5\%高かった」）．

\subsection{考察}
なぜそうなったのかを分析する（重要）．

「なぜ提案手法はうまくいった？」「逆にうまくいかなかったケースはある？それはなぜ？」など．これがないと実験レポート止まりになる．
% --- 結論 ---
\section{結論}
この論文の結論およびまとめと，今後の展望とかを書く．
アブストの再掲にならないよう，より具体的な結果を交えてまとめる．
本研究で解決しきれなかった課題や新しく見つかった問題点を正直に書く．「今後は〇〇への拡張が期待される」などで結ぶ．

結論まで書き終わった後に序論に戻ると，最初に宣言したことと最終結果にズレが生じていることがある．整合性を取るため，序論は最後に見直して修正する．

% --- 謝辞 ---
\section*{謝辞}
ここに謝辞を書く．



% --- 参考文献 ---
\begin{thebibliography}{99}
    \bibitem{1}
    Feng, R. and Shimizu, Y. (2013). On a generalization from ruin to default in a L\'{e}vy
insurance risk model. \textit{Methodology and Computing in Applied Probability}, \textbf{15}, (4), 773–802.

    \bibitem{2}
    Long, H.; Shimizu, Y. and Sun, W. (2013). Least squares estimators for discretely
    observed stochastic processes driven by small L\'{e}vy noises. \textit{Journal of Multivariate Analysis},
    \textbf{116}, 422–439.
    
    \bibitem{3}
    清水泰隆 (2019). 統計学への確率論，その先へ． 内田老鶴圃．    
\end{thebibliography}

% ---ここから付録---
\appendix % これ以降,節番号が A, B... に変わる
\section{本テンプレート使用方法}
\label{chap:manual}

\subsection{参照の自動化（cleveref）}
図表や式の番号を参照する際，\verb|\ref| は単に「1.1」のような数字を返すだけ．
そのため書き手が手動で「図\verb|\ref|」「式(\verb|\ref|)」と入力する必要があった．
\texttt{cleveref} パッケージの \verb|\cref| を使用すると，参照対象の種類（図，表，式，定理など）を自動判定し，適切な接頭辞を付与してくれる．

\begin{itemize}
    \item 例: \verb|\cref{eq:sample}| $\to$ \cref{eq:sample}
\end{itemize}


\subsection{定理環境と証明}
定義，定理，補題，命題，系など節ごとに自動で番号付けされる（例：\cref{thm:sample}）．

\begin{definition}[定義の名前]
    ここに定義を書く．オプション引数（\texttt{[ ]}）にタイトルを入れる．
\end{definition}

\begin{theorem}[定理の名前]
    \label{thm:sample} % 参照するためのラベル付け
    ここに定理を書く．
\end{theorem}

\begin{corollary}
    ここに系を書く．
\end{corollary}

\begin{proof}
    証明環境も定義済み．証明終了時にQ.E.D.が右端に配置．
    \[ 1+1=2 \]
\end{proof} 


\subsection{数式ショートカット}

\begin{table}[H]
    \centering
    \caption{定義済みマクロ一覧}
    \begin{tabular}{lll}
        \toprule
        コマンド & 出力 & 意味 \\
        \midrule
        \verb|\E| & $\E[X]$ & 期待値 \\
        \verb|\pb| & $\pb(A)$ & 確率\\
        \verb|\var| & $\var(X)$ & 分散 \\
        \verb|\cov| & $\cov(X,Y)$ & 共分散 \\
        \verb|\ind| & $\ind_{A}$ & 指示関数 \\
        \verb|\R|, \verb|\N| & $\R, \N$ & 実数,自然数 \\
        \verb|\pto| & $X_n \pto X$ & 確率収束 \\
        \verb|\dto| & $X_n \dto X$ & 分布収束 \\
        \verb|\asto| & $X_n \asto X$ & 概収束 \\
        \verb|\iid| & $X_i \iid F$ & 独立同分布 \\
        \verb|\Nd| & $X \sim \Nd (0,1)$ & 正規分布 \\
        \verb|\pdif{f}{x}| & $\pdif{f}{x}$ & 偏微分 \\
        \verb|\argmax_{a \in A}| & $\displaystyle\argmax_{a \in A}$ & argmax (argminも同様)\\
        \bottomrule
    \end{tabular}
\end{table}


\subsection{数式番号の確認}
\begin{align}
    y &= x^2 \label{eq:sample}
\end{align}
この式は \cref{eq:sample} である．

\subsection{TODOメモの活用}
執筆中，後で直したり先生に確認する箇所は \texttt{todonotes} パッケージが便利．
\todo{メモ内容}
コマンドは \verb|\todo{メモ内容}|．
提出時はパッケージオプションに \texttt{disable} を指定すると一括非表示になる．


\section{LaTeX 環境構築}
LaTeXを動かす環境として，自分のPCにインストールして使うローカル環境と，Cloud LaTeX / OverleafなどのWebアプリ環境がある．

個人的には圧倒的にローカル環境（VS Code）を強く推奨．
\begin{itemize}
    \item 爆速ビルド: クラウドの待ち時間（数秒〜数十秒）がなく，保存するたびに一瞬でプレビューが更新される．論文修正の試行回数が多い終盤に効いてくる．
    \item エディタ: VS Codeがそもそも優秀．予測補完や色分けによる視認性などが神．
    \item Git完全対応: 論文のバージョン管理ができる．「昨日の状態に戻したい」が瞬時に可能．
    \item AI支援: GitHub Copilotが使える．コードを勝手に補完してくれるので，執筆速度が上がる．
    \item 容量制限なし: 画像を大量に貼っても容量オーバーにならない．
\end{itemize}

WindowsはTeX Live，macOSはMacTeXをインストール．特にMacはすぐ終わる． \\

\href{https://code.visualstudio.com/download}{VS Codeインストールはこちら} \\

\href{https://texwiki.texjp.org/?TeX%E5%85%A5%E6%89%8B%E6%B3%95}{TeXインストールはこちら}
% ---ここまで付録---


\end{document}